% Standalone problem document, generated from the adjacent README.md.
% Compile with XeLaTeX. Mathematical content is maintained in Markdown.
\documentclass[11pt,a4paper]{article}
\usepackage[fontsize=10.5pt]{scrextend}
\usepackage[margin=22mm,headheight=15pt,headsep=9mm,footskip=11mm]{geometry}
\usepackage{fontspec}
\setmainfont{texgyrepagella}[Extension=.otf,UprightFont=*-regular,BoldFont=*-bold,ItalicFont=*-italic,BoldItalicFont=*-bolditalic]
\setsansfont{texgyreheros}[Extension=.otf,UprightFont=*-regular,BoldFont=*-bold,ItalicFont=*-italic,BoldItalicFont=*-bolditalic]
\setmonofont{texgyrecursor}[Extension=.otf,UprightFont=*-regular,BoldFont=*-bold,ItalicFont=*-italic,BoldItalicFont=*-bolditalic,Scale=MatchLowercase]
\usepackage{amsmath,amssymb}
\usepackage{unicode-math}
\setmathfont{texgyrepagella-math.otf}
\usepackage{xcolor,microtype,enumitem,needspace,fancyhdr}
\usepackage{bookmark}
\definecolor{ink}{HTML}{17344B}
\definecolor{accent}{HTML}{197D80}
\definecolor{muted}{HTML}{596773}
\hypersetup{colorlinks=true,linkcolor=accent,urlcolor=accent,pdftitle={FR-08 - Hadamard
matrices at every admissible
order},pdfauthor={Open Problems in Numerical Linear Algebra}}
\urlstyle{same}
\setlength{\parindent}{0pt}
\setlength{\parskip}{5pt plus 1pt minus 1pt}
\linespread{1.04}
\setlength{\emergencystretch}{3em}
\widowpenalty=10000
\clubpenalty=10000
\displaywidowpenalty=10000
\allowdisplaybreaks[1]
\setlist{nosep,leftmargin=1.5em}
\providecommand{\tightlist}{\setlength{\itemsep}{0pt}\setlength{\parskip}{0pt}}
\providecommand{\pandocbounded}[1]{#1}
\pagestyle{fancy}
\fancyhf{}
\fancyhead[L]{\sffamily\footnotesize\color{muted}OPEN PROBLEMS IN NUMERICAL LINEAR ALGEBRA}
\fancyhead[R]{\sffamily\bfseries\footnotesize\color{accent}FR-08}
\fancyfoot[L]{\parbox[b]{0.9\textwidth}{\sffamily\fontsize{8}{10}\selectfont\color{muted}Editorial ratings; a bounded search cannot certify that a problem remains unsolved.\\Literature check: 2026-09-10}}
\fancyfoot[R]{\sffamily\footnotesize\color{muted}\thepage}
\renewcommand{\headrulewidth}{0.3pt}
\renewcommand{\headrule}{\hbox to\headwidth{\color{accent}\leaders\hrule height \headrulewidth\hfill}}
\makeatletter
\renewcommand{\section}{\Needspace{5\baselineskip}\@startsection{section}{1}{0pt}{12pt plus 2pt minus 2pt}{4pt}{\sffamily\bfseries\large\color{ink}}}
\renewcommand{\subsection}{\Needspace{4\baselineskip}\@startsection{subsection}{2}{0pt}{9pt plus 1pt minus 1pt}{3pt}{\sffamily\bfseries\normalsize\color{ink}}}
\makeatother
\setcounter{secnumdepth}{0}
\begin{document}
{\sffamily\small\color{muted}Frames and matrix designs\par}
\vspace{3pt}
{\raggedright\hyphenpenalty=10000\exhyphenpenalty=10000\sffamily\bfseries\fontsize{21}{24}\selectfont\color{ink}Hadamard
matrices at every admissible order\par}
\vspace{7pt}
{\sffamily\small\textbf{Difficulty:} extreme\quad\textbf{Importance:} broadly
interesting\par}
{\sffamily\small\bfseries\color{accent}Status: Partially resolved\par}
\vspace{4pt}
{\color{accent}\hrule height 0.8pt}
\vspace{6pt}
\textbf{Rating rationale:} This longstanding all-order sign-design
conjecture underlies optimal real transforms and has broad consequences
in combinatorics, coding and experiment design.

Is it true that, for every positive integer \(m\), there exists a matrix
\[
H\in\{-1,1\}^{4m\times4m}
\quad\text{such that}\quad H^{T}H=4mI_{4m}?
\] Such an \(H\) is a real Hadamard matrix. The question concerns all
admissible orders, rather than a particular currently unconstructed
order.

\subsection{Relevance}\label{relevance}

The normalized matrix \(H/\sqrt{4m}\) is an orthogonal transform with
entries of equal magnitude. Thus this is an existence problem for
perfectly conditioned, maximally flat real sign transforms. It differs
from Cryer's complete-pivoting growth conjecture already in the
collection, which asks about elimination on matrices that already exist.

\subsection{References}\label{references}

\begin{itemize}
\tightlist
\item
  A. Bahmanian and S. Suda, \emph{Hadamard Hypercubes}, arXiv:2605.16722
  (2026), §1, opening definition and Conjecture 1.1
  (\href{https://arxiv.org/html/2605.16722v1}{primary manuscript}).
\item
  A. F. Ramos, D. B. Hulak, and R. J. G. B. de Queiroz, \emph{Multiplier
  obstructions for Legendre pairs of length 333}, arXiv:2607.20765
  (2026), abstract and introduction; a restricted construction problem,
  not a general obstruction
  (\href{https://arxiv.org/abs/2607.20765}{primary manuscript}).
\end{itemize}

\subsection{Status check --- 2026-09-10}\label{status-check-2026-09-10}

The May 2026 source states the universal conjecture explicitly. Searches
for ``Hadamard conjecture'', ``proof'', ``2026'', and recent
construction papers found no general resolution. Individual-order
construction announcements do not settle this statement. Because such
announcements can rapidly change the smallest unknown order, no
particular order is asserted here to remain open. This bounded
literature search is not a proof of open status.

\textbf{Audit update (2026-09-10):} Rechecked the May 2026 Conjecture
1.1 and searched for later Hadamard constructions and general proofs.
Established infinite construction families do not cover every multiple
of four; no newly verified general solution was located. This is a
bounded literature check, not a proof that no solution exists.
\end{document}
